\documentclass[pdflatex,sn-basic]{sn-jnl}
\usepackage[T1]{fontenc}
\usepackage[utf8]{inputenc}
\usepackage{amsmath,amssymb,graphicx}
\usepackage{microtype}
\hypersetup{hidelinks}
\setlength{\emergencystretch}{2em}

\begin{document}
\title[Terminal rescheduling and stop holding]{Terminal rescheduling and stop holding with a fixed bus fleet: passenger gains and circulation risks}
\author{\fnm{Author details} \sur{provided separately}}
\abstract{%
Moving departures toward demand can reduce waiting without adding trips, but the same buses must return in time to serve the opposite direction. We study whether stop holding adds value to terminal rescheduling when both decisions draw on the same vehicle time. FreqDuet uses an upper policy to revise unreleased departures and a lower policy to hold at stops, with holding bounded by projected turnaround availability. In a simulation informed by data for Boston's Route 111, four frozen policy pairs were evaluated on 16 new paired demand days with 14 buses and 262 trips. Compared with a dynamic rule using the same action set and trip budget, the learned controller reduced journey exposure by 4.26 seconds per passenger (95\% paired interval: 0.71 to 7.74 seconds saved), improved arrival regularity, and reduced readiness delay. Relative to the same learned architecture using unfiltered passenger-count inputs, harmonic policy inputs reduced journeys by 4.56 seconds per passenger, but did not meet the joint regularity and readiness criterion. Both learned variants retained the same harmonic input to the budget ledger. \input{generated/abstract_extension} The results concern headway-managed service with retiming, not adherence to a published timetable.
}
\keywords{Bus operations, terminal dispatch, stop holding, vehicle circulation, headway management}
\maketitle

\section{Introduction}

A bus can spend the same minute in three different ways: serving customers sooner, holding to regulate spacing, or returning to cover the next scheduled trip. Allocating more trips to a peak helps only if vehicles can actually depart. Likewise, a hold that improves headways at one stop can make the bus unavailable when the opposite terminal needs it. We ask whether joint dispatch and holding improve passenger service under a fixed physical fleet and trip roster, and how much holding adds once a learned dispatch policy is already in place.

Gap-feedback holding can moderate bunching but incurs time for riders already aboard \citep{daganzo2009headway,xuan2011dynamic,berrebi2018comparing}. Optimization-based methods account for capacity, waiting and holding costs \citep{eberlein2001holding,hickman2001analytic,delgado2012much,sanchez2016realtime}. Dispatch research has studied how to shift a fixed service supply toward higher demand \citep{newell1971dispatching}. \citet{sanchez2017schedulefree} examine real-time headway-managed dispatch and stop-skipping under operator resource constraints. Most directly, \citet{gkiotsalitis2021rescheduling} jointly optimizes departure rescheduling and bus holding with explicit original-timetable constraints, evaluating excess waiting against in-vehicle time in an operator-data simulation. Our emphasis is different: in a closed physical fleet, we measure what the same dispatch policy gains or loses when changes to holding alter the pool of buses available for subsequent departures. A same-authority rolling comparator extends the comparison beyond fixed timetables.

This paper contributes (i) an executable terminal-to-stop coordination rule that caps discretionary holding with a trip budget, (ii) paired comparisons that separate the effect of the complete learned controller, the choice of actor inputs, and the use of stop holding under a common upper policy, and (iii) a vehicle-level account of a circulation failure. The point is not to show that frequency separation guarantees better operations. It is to establish which observed gains survive a realistic base comparison and which disappear when a component is changed.

\section{Operating problem and methods}

\subsection{Service concept}
The service is a bidirectional, headway-managed route for passengers who arrive without targeting a published departure or a timed transfer. The original nominal roster provides 131 trips per direction at a 360-second nominal spacing. Vehicles come from a single fixed physical fleet, reversing direction at each terminal. Every 900 seconds the planner updates the unreleased departure plan; a trip already released, queried for dispatch, or eligible for a bus is not moved. The projection assigns an active block of gaps and a compensating tail while retaining the trip roster and final departure no earlier than the demand cutoff. Dispatch also requires a ready vehicle and a 120-second minimum actual spacing. Neither that floor nor the planned spacing guarantees a particular realized headway distribution.

The model commits trips at the simulator decision boundary. It does not promise a public notice period and does not model a passenger reaction to announcements. In operation, a headway-managed route would need rolling passenger information and an agency-defined lock for already advertised departures. For timetable-based service or guaranteed connections, early departures would need a separate constraint: a plan that runs with little readiness delay may still be unsuitable. The untimed-arrival assumption is most plausible in frequent periods. The study does not assume it remains valid during low-frequency shoulders or near last departures; demand generation stops at 19:03, and the roster and evaluation include a separate clearance period.

\subsection{Controller and observations}
The upper policy selects among 15 actions every 15 minutes. Each is a pair of a quota-adjusted spacing multiplier in $\{0.78,0.89,1.00,1.12,1.24\}$ and a directional tilt in $\{-36,0,36\}$ seconds. The inventory guard attenuates the tilt when directional ready stock is low. The same projection is used for all corrected comparisons. Along the route, decisions are made after boarding and alighting using the state of the bus and causal estimates of the adjacent gaps. Stop holding is restricted to the bins $\{0,5,10,15,20,25\}$ seconds.

The turnaround ledger predicts ready times using current terminal stock, active vehicle position, residual stop time, and a recent direction-specific median nonholding trip time. At dispatch and at least every 60 seconds, it sorts predictions at each terminal and matches them to unlaunched trips. It deducts a 20-second planning buffer from predicted slack and converts the positive remainder to a budget with retention 0.50. The budget is capped at 85 seconds; stronger opposite-direction demand and less than 75 seconds of slack instead trigger a 0.40 multiplier on the allowance and a 30-second cap. An initially unmatched trip gets 37.5 seconds. A later calculation can shrink but not replenish the budget, and actual holding is debited. This is a hard action limit, not a guarantee that a predicted vehicle will return on time.

The causal boarding proxy divides observed boardings by the observed arrival gap (with fallback, clipping, and a fullness adjustment) and updates a held stop rate. A 20-second sample contributes one third of its per-minute rate to the closed one-minute bin. A four-harmonic recursive least-squares estimator with a static demand prior separates the smooth component and residual. The aligned policy sends directional low-frequency summaries to the upper actor and signed local residual summaries to the lower actor. Both use shared physical features. The upper state contains 13 physical and clock features and eight demand features; the lower state has 15 physical and clock features, including the mask limit, and four demand features. Neither actor reads the station's exact post-boarding queue. Swapped uses matched summaries of the opposite band. NoFreq uses the same summaries of the completed, unfiltered APC bin and its one-bin change for both actors. The ledger still reads the same harmonic low-frequency signals in these variants; the comparisons change actor inputs, not all harmonic information in the system.

Discrete entropy-regularized actor--critic policies \citep{haarnoja2018soft,christodoulou2019discrete} are trained with the same critic and update budget. The upper reward penalizes simulated waiting, exposure to discretionary holding, an excess-waiting proxy, and actual vehicle-readiness delay; it also includes a remaining-quota potential. The lower reward subtracts onboard delay and predicted slack penalties from a two-gap waiting proxy. With observed forward and predicted follower gaps $g_f,g_b$, targets $H_f,H_b$, and a downstream boarding-rate proxy $\lambda$, the two-sided term is
\begin{equation}
G(u)=\lambda\frac{H_f+H_b}{2}u
  \left[\frac{g_b}{H_b}-\frac{g_f}{H_f}
    -\frac{u}{2}\left(\frac1{H_f}+\frac1{H_b}\right)\right].
\end{equation}
At training time the downstream rates are reconstructed from subsequent boardings on that trip. This changes a reward label, not a runtime observation. The lower critic uses a zero-hold action reference and a Lagrange penalty for regularity overshoot, onboard load, and tight projected slack; neither the proxy nor the cost critic ensures a non-harmful action in operation.

\subsection{Benchmarks and evaluation}
The corrected case uses Fall 2025 seasonal automatic passenger count averages and stop geometry for MBTA Route 111. It is a 23-stop data-informed simulation: the published trips are not replayed, origin--destination flows are reconstructed, running speeds are constructed, and control uses a synthetic timetable. Exogenous traffic and demand are paired by scenario tape. V503 swaps the actor frequency inputs on 16 held-out days. V504 evaluates four 200-episode seeds over a different 16-day test set against both matching-budget, matching-architecture NoFreq and a dynamic rule (Manual). Manual picks from the same 15 release actions by nearest directional square-root target based on the raw completed boarding bin, and uses a causal, load-aware hold scorer based on the same waiting and slack proxy. It runs with the same six holding bins, projection, and budget. Shared constraints do not imply matched actual holding dose.

To attribute holding separately, a closed-loop intervention keeps each Aligned upper actor fixed and substitutes either zero discretionary holding or the Manual hold scorer for its lower policy. The upper actor replans from that arm's own realized vehicle stock and holding history. Simply replaying an identical departure plan would suppress this feedback. The intervention reuses the same 16 V504 days, so it is a mechanism comparison, not a new external confirmation. The additional Rolling baseline combines the same Manual hold scorer with an upper search over the native 15 projected plans. It forecasts the next eight trips per direction, recycles the predicted returning fleet, and scores squared-gap expected waiting plus readiness cost (one nominal gap's arrivals per second of delay). Its inputs are the same causal stock, position, current stop time, recent nonholding times, and completed APC bins. It is not a reproduction of the integer nonlinear program of \citet{gkiotsalitis2021rescheduling}.

We report waiting plus onboard exposure per generated passenger, actual arrival-headway CV$^2$, and actual vehicle-readiness delay from first dispatch eligibility under the active plan. Exposure includes passengers still waiting or onboard at the fixed end of evaluation. We also report the absolute actual departure displacement from the original nominal timetable, decomposed into positive earliness and lateness, and note whether every scheduled trip and passenger was served. The common original timetable is an operating reference, not a passenger promise in this simulation. Paired bootstraps resample training seeds and demand days independently (30,000 draws), with deterministic rule scores shared rather than treated as additional trained seeds. The joint criterion requires the upper 95\% endpoint to be negative for journey and nonpositive for CV$^2$, readiness and absolute deviation (V503 uses journey, CV$^2$ and readiness).

\section{Results}

\subsection{Corrected fixed-fleet comparison}
Table~\ref{tab:main} gives the main results. All 144 evaluation paths operated the complete roster with the fixed fleet and zero unserved passengers. Aligned reduced journeys relative to Manual in all four seed means and 13 of 16 day means. The reduction was 0.38\% of the Manual journey mean, including 3.93 seconds less waiting and 0.33 seconds less onboard time; the onboard interval included zero. Multiplying the paired differences by each scenario's generated passenger count gives \input{generated/person_hours} passenger-hours saved per operating day. With the same budget constraints, the policies still held very differently: 3.46 versus 11.73 seconds per trip. The result supports the combined learned controller on this corridor, not an improvement per unit of holding.

\begin{table}[!ht]
\centering\small
\caption{Main paired comparison on 16 new MBTA-informed simulation days (V504). Entries are Aligned minus the comparator, with pointwise 95\% seed-by-day bootstrap intervals. Negative values favor Aligned. Absolute offset is measured against the original nominal timetable.}
\label{tab:main}
\input{generated/v504_table}
\end{table}

Relative to NoFreq, harmonic actor inputs reduced journeys on 15 of 16 day means and in three of four seed means. The point estimate for arrival regularity also favored Aligned, but its interval did not establish non-harm. Average readiness delay was higher. An exploratory NoFreq--Manual comparison is useful for interpreting the large readiness gain: even without harmonic actor inputs, NoFreq lowered readiness delay by 98.35 seconds per trip relative to Manual, while the journey difference was only $+0.30$ seconds per passenger (interval $[-5.84,+5.99]$). The large readiness advantage against Manual is therefore not evidence of a harmonic actor effect. V503, the matched band-swap study on a separate day set, also failed its joint criterion; the journey difference was $-4.28$ seconds per passenger (interval $[-10.56,+1.45]$).

\subsection{Holding attribution and rolling benchmark}
\input{generated/holding_results}

\subsection{What readiness and timetable displacement measure}
The Aligned controller had 3.65 seconds of readiness delay per trip, measured from eligibility under its own active plan, and 515.42 seconds of absolute displacement from the original 360-second timetable. These are not contradictory: trips can move substantially on the calendar while a bus is usually ready to run the revised plan. Averaged over every trip (including zero contributions), Aligned had 319.70 seconds of earliness and 195.72 seconds of lateness, compared with 257.58 and 503.47 seconds with Manual and 743.01 and 356.11 seconds with NoFreq. Thus the lower absolute displacement against Manual was driven by less lateness, not uniformly by less earliness. An eight-minute mean absolute plan displacement is not an eight-minute change in experienced passenger waiting.

\subsection{A measured circulation failure in simulation}
Seed 50314 on day 53020003 incurred 150.50 seconds of readiness delay per trip, versus 5.34 seconds with NoFreq. The pair accounts for about 98\% of the overall mean readiness difference. The observation remains in all primary results. We examine it as an already identified failure, not as a randomly sampled operating day or an estimate of how often failures occur.
\input{generated/case_results}

\subsection{Independent running-time and terminal check}
\input{generated/field_results}

\subsection{Relationship to the earlier implementation}
The archived V500--V502 studies used a nominal 60-second conversion at each 20-second APC update, and the old actor read the exact simulated post-boarding queue. The full old controller reduced mean journeys by 14.32 and 14.02 seconds per passenger against fixed dispatch with no holding on the MBTA-informed and synthetic tidal corridors, respectively. Only the first interval excluded zero when training seeds were included. The archived rolling rule was faster than the old learned controller by 5.20 and 16.32 seconds per passenger, respectively, but had higher readiness delay. In the old eight-day non-learning development study, adding constrained holding to the same dispatch rule increased mean journeys by 3.37 and 3.07 seconds per passenger. These are historical context, not tests of the current controller or estimates of retiming in isolation. The corrected variant changed the upper action set, fixed retention, and matched the feature comparisons; subtracting the old 9.59-second harmonic-input estimate from the new 4.56-second estimate would not isolate a unit-conversion effect.

\section{Discussion and conclusions}

The corrected study shows why better passenger time and operational reliability need separate evaluation. Against the specified dynamic rule, the learned controller had both lower journey exposure and more regular, less delayed operations; against the matched raw-input learner, journeys fell but the joint criterion was not met. Consequently, the gain from the full controller is better supported than a general gain from assigning different frequency bands to its two policies. This distinction matters when different features also change the realized timetable and use of the holding budget.

\input{generated/holding_discussion}

\subsection{Limitations}
The corrected evaluation uses one data-informed route model and four training seeds. Reusing paired days in the mechanism comparison does not create another independent confirmation. Boardings under capacity constraints are a censored demand proxy, and projected trip times can misrepresent future vehicle stock. The field observations benchmark running times and selected same-vehicle turnaround gaps; they do not validate the simulated controller under deployment. In published-timetable service, a controller should not depart early against an advertised promise. The results do not measure timed-transfer outcomes, off-peak schedule-sensitive rider behavior, or the consequences of adding a passenger notice or departure-advance constraint.

\appendix
\section{Implementation details}

\subsection{Actor inputs}
The upper operational block uses phase sine and
cosine, remaining day fraction, each directional remaining-trip fraction
relative to the remaining day, quota-available gaps, directional ready
counts, total active fleet fraction, the cycle-based spacing floor, recent
budget utilization and preceding interval readiness delay. For each
direction, the four demand statistics are the sum, scaled spatial root
mean square, scaled spatial standard deviation and aggregate one-bin
change. For the lower actor the operational inputs are the dynamic action
limit, remaining holding, holding per remaining stop, estimated turnaround
slack, overdue trips at the next terminal, two relative headways, two-sided
availability, analytical gap-based hold target, the two planned target
headways, onboard fraction, route progress, and phase sine and cosine.
Its demand statistics are the local quantity and the mean of the next
two stops, each with its one-bin change. This summary and its normalization
are shared by the corrected input comparisons.

\subsection{Updates}
Each policy and critic has two 96-unit hidden layers.
The reward critics use six members and a pessimistic mean-minus-standard
deviation target. The upper replay uses six-step discounted returns
($\gamma=0.996$), a capacity of 20,000 and batches of 48. The lower replay
uses one-step returns ($\gamma=0.92$), a capacity of 180,000 and batches of
128. Both learning rates are $3\times10^{-4}$; the lower multiplier uses
$10^{-3}$ and a cost target of 0.020. The four final seed pairs use the
first two days as replay warm-up and make 32 upper and 64 lower updates
on each subsequent training day. Final evaluation uses greedy action
selection with the same action mask as training.

\subsection{Limits and forecasts}
The dynamic lower mask multiplies the
normalized remaining action allowance by the 25-second bin maximum;
the zero bin is always legal. The upper spacing floor estimates the
two directional trip times plus 40 seconds divided by the fleet size,
then scales by 0.85, 0.90 or 0.97 as terminal stock falls from at least
three to two or fewer than two vehicles. The result is clipped to
$[250,335]$ seconds, with an upper requested gap of 455 seconds. The
operating projection falls back to uniform residual gaps when the
two-block allocation is infeasible. The minimum actual dispatch
spacing enforced by the simulator is a separate 120-second rule.


\section*{Statements and Declarations}
\noindent\textbf{Data and code availability.}
The accompanying review archive contains the frozen simulation workbooks, paired result summaries, and code to recreate the tables and figure. The associated study protocols and simulation code are archived with the project. Agency data are credited in the methods and the review archive; they are not included in the controller evaluation as realized futures.

\noindent\textbf{Generative AI.}
Google Gemini and OpenAI Codex assisted in revision of the text and implementation. The report uses the numerical outputs of the archived code and paired study records.

\bibliography{refs}
\end{document}
